\documentclass[12pt]{article}

\usepackage{amsmath}
\usepackage[margin=1in]{geometry}
\usepackage{fancyhdr}
\usepackage{amsthm}
\newtheorem{lemma}{Lemma}
\usepackage{bm}
\usepackage{braket}
\usepackage{amssymb}
\usepackage{tikz}
\usepackage{bbm}
\usetikzlibrary{shapes.geometric}
\usepackage{enumerate}
\usepackage[shortlabels]{enumitem}


\pagestyle{fancy}
\fancyhead[l]{Hudson Benites}
\fancyhead[c]{Day 2}
\fancyfoot[c]{\thepage}
\renewcommand{\headrulewidth}{0.2pt}
\setlength{\headheight}{15pt}
\fancyhead[r]{\today}

\begin{document}
\begin{enumerate}[label=\Alph*.]
    \item \begin{enumerate}[label=\theenumi\Alph*]
        \item We can see that $\bm{\alpha}$ is smooth because each of its component functions are polynomials, which are smooth. We show that $\bm{\alpha}$ is regular by showing that $\bm{\dot{\alpha}}(u)\neq 0$ for all $u$. We see that $\bm{\dot{\alpha}}=(1, 2u)$, which can never be the zero vector. We can then use $\bm{\dot{\alpha}}$ to compute the unit tangent vector field $\mathbf{t}$ (now known as $\mathbf{T}$). We compute
    \begin{equation*}
        \mathbf{t}(u)=\frac{\bm{\dot{\alpha}}}{\vert\bm{\dot{\alpha}}\vert}=\frac{\begin{bmatrix}
            1\\
            2u^2
        \end{bmatrix}}{\sqrt{1+4u^2}}.
    \end{equation*}
    \item We can see that $\bm{\alpha}$ is smooth because each of its component functions are polynomials, which are smooth. We show that $\alpha$ is not regular by showing that $\dot{\alpha}=0$ for some $u$. We see that $\dot{\alpha}(u)=(3u^2, 6u^5)$, which is the zero vector when $u=0$. Therefore, because $\alpha$ is not regular,
    it is not possible to compute the unit tangent vector field $\mathbf{t}$ because it is not defined when $\dot{\alpha}=\mathbf{0}$.
\end{enumerate}
    \item \begin{enumerate}[label=\theenumi\Alph*]
        \item The speed of $\bm{\alpha}$ is given by $\sqrt{a^2\cos^2(u)+b^2\sin^2(u)}$. To find the minima and maxima of the speed, we can instead optimize the square of the speed. We see that ${\vert\bm{\dot{\alpha}}\vert}^2=a^2\cos^2(u)+b^2\sin^2(u)$. We differentiate to find the minima and maxima. We compute that \begin{equation*}
            \frac{d}{du}{\vert\bm{\dot{\alpha}}\vert}^2=2a^2\sin(u)cos(u)-2b^2\sin(u)\cos(u)=2\sin(u)cos(u)(a^2-b^2).
        \end{equation*} By observation we see that this equals 0 whenever $a=b$ (which never happens because $a>b>0$), and when $\sin(u)=0$ or $\cos(u)=0$. This occurs $\frac{\pi}{2}k$ for $k\in \mathbb{Z}$, making $\sin$ and $\cos$ equal to 0 at alternating $k$. At $\pi k$, the speed is \begin{equation*}
            \vert\bm{\dot{\alpha}}\vert = \sqrt{a^2\cos^2(\pi k)} = a.
        \end{equation*} Similarly, at $\pi k+ \frac{\pi}{2}$, the speed is \begin{equation*}
            \vert\bm{\dot{\alpha}}\vert = \sqrt{b^2\sin^2(\pi k + \frac{\pi}{2})}=b.
        \end{equation*} By assumption, $a>b$, and therefore the speed is maximized at all points $\pi k$ and minimized at all points $\pi k + \frac{\pi}{2}$ for $k \in \mathbb{Z}$.
        \item No. The parameterization $u$ does not need to be related to $t$, and therefore B.A only tells us the speed relative the $u$, not the actual time $t$.
    \end{enumerate}
    \item \begin{enumerate}[label=\theenumi\Alph*]
        \item We compute the angle $\theta$ between $\bm{\alpha}$ and $\bm{\dot{\alpha}}$ by computing $\arccos(\frac{\bm{\alpha} \cdot \bm{\dot{\alpha}}}{\vert \bm{\alpha}\vert \vert \bm{\dot{\alpha}} \vert})$. Let $a=ce^{\lambda u}\cos(u)$ and $b=ce^{\lambda}\sin(u)$, noting that $a^2+b^2=c^2e^{2\lambda u}$. Then $\bm{\alpha}=\begin{bmatrix}
        a\\
        b
        \end{bmatrix}$ and
        $\bm{\dot{\alpha}}=\begin{bmatrix}
            \lambda a - b\\
            \lambda b + a
        \end{bmatrix}$. We compute that $\vert\bm{\alpha}\vert=\sqrt{a^2+b^2}=ce^{\lambda u}$. We compute that $\vert \bm{\dot{\alpha}} \vert=\sqrt{(\lambda a -b)^2 + (\lambda b + a)^2}$. We see that the radicand simplifies to
        $\lambda^2a^2 + \lambda^2b^2 + a^2 + b^2$, which is equivalent to $(a^2+b^2)(\lambda^2 + 1)$. Plugging back in, we get that \begin{equation*}
            \vert \bm{\dot{\alpha}} \vert = \sqrt{c^2e^{2\lambda u}(\lambda^2 + 1)}=ce^{\lambda u}\sqrt{\lambda^2 + 1}.
        \end{equation*} We also compute $\bm{\alpha}\cdot \bm{\dot{\alpha}}$ elementwise. We see that \begin{equation*}
            \bm{\alpha}\cdot\bm{\dot{\alpha}} = \begin{bmatrix}
            a & b
            \end{bmatrix}\begin{bmatrix}
                \lambda a - b\\
                \lambda b + a
            \end{bmatrix}=a(\lambda a - b) + b(\lambda b + b)=\lambda c^2e^{2\lambda u}.
        \end{equation*}
        Putting it all together, we see that
        \begin{equation*}
            \arccos(\frac{\bm{\alpha} \cdot \bm{\dot{\alpha}}}{\vert \bm{\alpha}\vert \vert \bm{\dot{\alpha}} \vert})=\arccos{\frac{\lambda c^2e^{2\lambda u}}{(ce^{\lambda u})(ce^{\lambda u}\sqrt{\lambda^2 + 1})}}=\arccos({\frac{\lambda}{\sqrt{\lambda^2 + 1}}}),
        \end{equation*} which is $\theta$ and constant.
        \item Recall that the arc length of a parameterized curve is given as a function of $u$ for a fixed $u_0$ by $s(u)=\int_{u_0}^{u}\vert \bm{\dot{\alpha(v)}}\vert dv$. Fix $u_0=0$. For this particular $\bm{\alpha}$, this amounts to\begin{equation*}
            s(u)=c\sqrt{\lambda^2 + 1}\int_{u_0}^{u}e^{\lambda v} dv=\frac{c\sqrt{\lambda^2 + 1}}{\lambda}(e^{\lambda u}-1).
        \end{equation*} We can solve for $u$ to find the inverse $u(s)$. We compute that
        \begin{equation*}
            u(s)=\frac{1}{\lambda}\ln(\frac{\lambda s}{c\sqrt{\lambda^2 + 1}}+1).
        \end{equation*}
        Using $u(s)$ to parameterize $\bm{\alpha(u)}$, we see that
        \begin{equation*}
            \bm{\alpha(u(s))}=c(\frac{\lambda s}{c\sqrt{\lambda^2 + 1}}+1)\begin{bmatrix}
                \cos (\frac{1}{\lambda}\ln(\frac{\lambda s}{c\sqrt{\lambda^2 + 1}}+1))\\
                \sin (\frac{1}{\lambda}\ln(\frac{\lambda s}{c\sqrt{\lambda^2 + 1}}+1))
            \end{bmatrix}.
        \end{equation*}
    \end{enumerate}
\end{enumerate}
\end{document}